\section{Result}

\subsection{Verbal and Interaction Strategies for Controlling Robot Motion}
For \textbf{RQ1}, we seek to characterize the verbal and interaction strategies participants use when controlling a robot's physical motion through natural language. Across the two tasks, participants varied not only in how they described motion, but also in how they controlled its temporal progression, established spatial reference frames, and reasoned about the robot's behavior. We organized these behaviors into three dimensions: movement command specification, command sequencing, and frame-of-reference management.

\subsubsection{Movement Command Specification}
Participants drew on five strategies to specify robot motion, often shifting between or combining them during interaction: quantitative measures, qualitative expressions, user-defined units, relative adjustments, and environmental references. \textbf{Quantitative measures} involved specifying movement using explicit magnitude such as centimeters, degrees, or time (P3, P4, P5, P8, P10). Participants gave commands such as \emph{``move forward 50 centimeters''} and \emph{``turn left 30 degrees.''} Time was also used as a proxy for distance, as in \emph{``walk straight for 5 seconds.''} In contrast, \textbf{qualitative expressions} relied on vague magnitude terms such as \emph{``a little bit,''} \emph{``a larger step,''} or \emph{``a really long time''} (P6, P7, P8, P9, P12). These expressions were often refined incrementally after participants observed the robot's movement, for example by asking it to move \emph{``a little bit more.''}

Some participants established their own movement scales through \textbf{user-defined units} (P5, P7, P9, P11). P5 explicitly defined a particular movement amount as \emph{``one unit,''} while others used measures such as \emph{``steps''} or wheel rotations without formally specifying their physical magnitude. Participants also used \textbf{relative adjustments} to specify new movements in relation to previous ones (P4, P5, P8, P9, P10, P12). Rather than providing a new absolute magnitude, they requested the same amount, \emph{``twice the amount,''} \emph{``5 more degrees,''} or simply instructed the robot to \emph{``keep going.''} Finally, we observed one participants used \textbf{environmental references} to describe desired movement in relation to objects or destinations rather than low-level motion parameters. For example, P1 instructed the robot to \emph{``go in between a frame and a cup.''} Because the robot could not visually ground these environmental references, participants quickly shifted away from this strategy toward primitive movement commands.

\subsubsection{Movement Command Sequencing}
Participants used three strategies to organize commands and regulate how robot motion unfolded over time: discrete step-and-stop control, user-terminated continuous commands (UTCs), and command batching. In \textbf{discrete step-and-stop control}, participants (P1, P2, P3, P6) issued short movement primitives, observed the result, and stopped the robot before continuing. Typical sequences included \emph{``Turn right... Stop''} and \emph{``Forward... Stop.''} Stop commands often marked either a movement boundary or the detection of an error, and were sometimes accompanied by repair expressions such as \emph{``No! Turn left''} or \emph{``Sorry, stop.''}

In \textbf{user-terminated continuous commands (UTCs)}, participants allowed the robot to continue moving until they later issued a stop. Unlike step-and-stop control, the terminal stop was generally anticipated as part of the control strategy rather than introduced primarily as a correction. Some participants (P1, P2, P4, P8, P9, P11) explicitly established this stopping arrangement in advance, as in \emph{``Keep moving forward until I say stop.''} Others (P9, P10, P11, P12) sustained the ongoing movement through repeated continuation cues, such as \emph{``Keep going. Keep going... Stop.''} Finally, in \textbf{command batching}, participants (P3, P5, P7, P8, P12) combined multiple movement primitives within a single utterance. Rather than issuing each instruction separately, they specified a sequence of actions, such as moving forward, rotating, and then moving forward again.

\subsubsection{Managing Frame-of-Reference Ambiguity}
Participants used several strategies to resolve ambiguity in directional commands such as left, right, forward, and backward, which could be interpreted from the participant's perspective, the robot's perspective, or another external frame of reference. Some participants \textbf{anchored directional commands} to the robot's body frame. P10 and P12 explicitly referenced the robot's perspective using expressions such as \emph{``Turn to your right a little bit''} or \emph{``Turn a little bit to your left.''} P11 instead used \textbf{orientation-independent terms}, such as \emph{``clockwise''} and \emph{``counterclockwise,''} to avoid left--right ambiguity. P11 also distinguished rotation from translation, using \emph{``spin''} to indicate rotation in place and \emph{``move''} or \emph{``go''} to indicate translation, as in \emph{``Stay in the same position; spin 45 degrees.''}

Other participants resolved uncertainty interactively. P6, P7, and P8 \textbf{asked about the robot's orientation} before continuing, using questions such as \emph{``Which way is your head?''}, \emph{``Which side is the front?''}, and \emph{``Wait, which way is forward?''} P7 described the ambiguity directly: \emph{``I don't know whether I should tell the direction according to my view or the robot's view.''} Participants also used \textbf{test-and-recalibrate} strategies when the reference frame remained unclear (P2, P8, P11). P8 tested forward and backward movements, while P11 confirmed the observed mapping with \emph{``Okay, so that's forwards.''} P2 instead reset the robot to a canonical orientation using \emph{``Turn to the front''} before continuing.

\subsection{Challenges and Interaction Needs}
For RQ2, we seek to learn about the challenges participants encountered as they adjusted and repaired verbal robot control, as well as the interaction support they expected across different task demands. Participants generally found the free navigation task more tolerant of imprecision because they could adapt the route as long as the robot reached the intended destination (P3, P8, P9, P10, P11), whereas the constrained path task following required finer-grained turns, closer trajectory adherence, and more frequent corrections (P5, P8, P11). Participants also reported mixed overall experiences: some found the interaction comfortable or manageable when the robot generally followed their instructions (P1, P4, P5, P7), while others described it as frustrating or tiring because of repeated corrections and the number of low-level commands required (P2, P3, P8, P11). These differences provide context for the challenges and interaction needs described below.

\subsubsection{Control and Execution Challenges}
The most common challenge was a mismatch between the motion participants specified and the motion the robot produced (P3, P4, P5, P7, P8, P9). This mismatch occurred even when participants used explicit numerical commands. Rotation was particularly difficult to predict (P3, P4, P9). P3 explained that \emph{``sometimes when I say maybe 30 degrees, it takes like 60 degrees or just like 90 degrees,''} while P9 observed that \emph{``every time the angle that it makes turns is kind of different.''} P4 similarly noted that \emph{``what I thought was not actually what the robot was responding to.''} Participants encountered similar problems with distance, speed, and straight-line movement (P5, P7, P8, P9). P5 reported that \emph{``when the robot moved forward, it was not exactly moving the same amount as I wanted it to.''} P7 had to \emph{``guess how many steps it takes''} and then adjust subsequent commands, while P8 repeatedly corrected the robot because it drifted slightly to the left or right rather than maintaining a straight trajectory.

Participants also struggled to translate spatial intentions into verbal quantities. P10 stated, \emph{``It is hard for me to estimate how long it should walk, as well as what exact degree it should rotate.''} As a result, participants sometimes relied on approximations, such as specifying movement duration rather than physical distance. These difficulties increased the interaction effort required to control the robot. P2 described the process as \emph{``tedious''} because every small movement required another instruction, while P8 recalled having to issue a forward command approximately 20 times, which was \emph{``tiring.''} Repeated micro-instructions therefore made verbal control burdensome, particularly when precise positioning required frequent observation and correction.

\subsubsection{Challenges in Diagnosing and Repairing Failures}
A further challenge was determining why the robot had behaved unexpectedly and how to recover from the resulting mismatch. Participants did not always know whether a failure originated in their command, the robot's interpretation, or its physical execution. Some attributed problems to the robot's interpretation of their language. P6, for example, stated, \emph{``I feel like maybe the numerics are not working.''} Others attributed discrepancies to physical execution. After observing an inaccurate turn, P11 reasoned, \emph{``It turned 170 because of friction, so let's try 190,''} and compensated by modifying the next command. Participants also considered whether their own wording had contributed to the problem. P10 reflected, \emph{``Maybe I'm not using the best command to communicate.''}

Because the source of failure was often uncertain, repair required participants to experiment with different responses. They reformulated commands, adjusted numerical parameters, or replaced an ongoing instruction with a new one. P4 described issuing a follow-up instruction to \emph{``break the old one and continue this new routine,''} while P5 preferred asking the robot to \emph{``repeat the last step.''} In some cases, participants also disregarded robot feedback when it conflicted with what they could directly observe. P1 instructed the robot to \emph{``Go forward, trust me,''} while P11 responded to an obstacle warning with \emph{``That's a cup, you don't have to go around.''} The effort required to diagnose and recover from these failures also shaped participants' expectations for how the robot should support the interaction.

\subsubsection{Interaction Needs}
Participants wanted the robot to provide enough feedback to make its interpretation and execution understandable. They expected it to acknowledge commands, state the action it intended to perform, and report completion or progress. P4 suggested brief acknowledgments such as \emph{``okay''} or \emph{``got it,''} while P9 wanted the robot to state, for example, \emph{``I am now going to turn right for 60 degrees,''} and confirm when the action was complete. Such feedback was particularly important for distinguishing whether an unexpected outcome resulted from misunderstanding or physical execution. Without feedback, P10 assumed that the robot \emph{``does not understand me,''} while P11 wanted to know whether a problem came from \emph{``a faulty robot''} or the \emph{``command.''} Participants also valued obstacle feedback (P2, P3, P7, P9), although its usefulness depended on what they could already observe. At the same time, they did not want excessive verbal communication. P1 stated that the robot should \emph{``talk back when necessary but shouldn't be talking back too much,''} and P6 suggested that routine information such as obstacle detection could instead be communicated through a non-verbal signal, such as a red light.

When commands were ambiguous, participants expected the robot to ask for clarification rather than silently select an interpretation. However, they also expected the robot to retain the resulting clarification. P8 noted that repeatedly asking what \emph{``a little bit''} meant would become \emph{``kind of annoying.''} Participants therefore favored clarification combined with memory, so that locally established meanings could be reused during the interaction. Participants further expected the robot to disclose its capabilities and limitations and to support calibration before or during the task. P4 suggested providing example command forms and a brief calibration procedure to reduce trial-and-error interaction. More generally, participants expected negotiation to remain proportional to the task. P7 stated that the robot should generally \emph{``strictly follow the user's instructions,''} while explaining its reasoning when it could not. P11 similarly emphasized that additional interaction effort should feel justified, stating, \emph{``I want my efforts to be necessary.''}